\section{The remaining reductions}\label{sec:exceptions}

In this section we deal with the reductions in \Cref{table:exceptions}, keeping the notation of \Cref{sec:local}. Recall that $\overline D=Z+C'$, that $C$ is the rational cuspidal divisor lifting $C'$, and $c=\deg C=d-\deg Z$. Most of the cases are dealt with by \Cref{lem:stab}, in which case we give the element $h$, the reason why $y=0$ (if we use it), and the relevant inequality for the gonality, which always comes from \Cref{prop:gonality}. All the statements that $h$ fixes a closed point of $X_1(p)_{\F_2}$ are \gitlink{degree19/stabilizers.m}{checked} on the elliptic curves representing these points. When we say that \textit{$D$ is a pullback from $X_H$}, we mean that $D=\pi^*D'$, where $\pi:X_1(p)\to X_H$ is the quotient map and $D'$ is an effective divisor on $X_H$ defined over $\Q$, with no cusps in its support. By \Cref{lem:orbits}, this is the case if $D$ is invariant under $H$ and $d<(p-1)/3$.

\subsection{The prime \texorpdfstring{$37$}{37}}

\begin{proposition}\label{prop:37}
We have $37\notin S(9)$ and $37\notin S(11)$.
\end{proposition}

\begin{proof}
We \gitlink{audit11/rational_two.m}{compute} that the common kernel of $A_3$ and $A_5$ on $H_1(X_1(37),\Z)\otimes\F_2$, which contains $J_1(37)(\Q)[2]$ by \Cref{lem:torsion} b), is $0$. So $y=0$ in \Cref{lem:stab}, and we have $\gon_\Q X_1(37)=18$. The exceptional reductions are the following \gitlink{audit11/stabilizers.m}{orbits}.

Suppose first that $Z=O_6$, one of $3$ closed points of degree $6$ fixed by $h=\dm8$, of order $6$. Then $c=d-6\in\{3,5\}$, and $18>d+c$, so $6\mid c$ by \Cref{lem:stab}, which is a contradiction.

Suppose now that $Z$ is one of the $2$ ordinary closed points of degree $9$, or the supersingular closed point of degree $9$, all of which are fixed by $h=\dm4$, of order $9$. Then $c\in\{0,2\}$ and $18>d+c$, so $hD=D$ and $9\mid c$. Hence $c=0$, $d=9$, and by \Cref{lem:orbits} the support of $D$ is a single orbit of $\dm4$, so the image of $D$ on $X_0(37)$ is a non-cuspidal rational point. The corresponding $j$-invariants are $-7\cdot11^3$ and $-7\cdot137^3\cdot2083^3$ \cite[Section 5]{MSD}, those of the curves \lmfdbec{1225}{b}{2} and \lmfdbec{1225}{b}{1}. As $E$ is a quadratic twist of one of these curves over $K$, the $x$-coordinate of $P$ is a root of the $37$-division polynomial of one of them which lies in $K$. But we \gitlink{degree09/P37/rational_isogeny_fibers.m}{compute} that the irreducible factors of these division polynomials have degrees $6,6,6,666$ and $18,222,222,222$, none of which divides $9$.

Suppose finally that $d=11$ and $Z=O_{10}$, so $c=1$. For one of the two orbits, $Z$ is fixed by $h=\dm 6$, of order $2$, and as $18>12$, \Cref{lem:stab} gives $2\mid 1$. The other orbit consists of the reductions of the closed points of degree $10$ in the supports of the trace $S$ of the point of order $37$ of \Cref{sec:lower} and of its images under $\Delta(37)$ \gitlink{audit11/decic37.m}{(we check this)}. Replacing $S$ by such an image we can assume $\overline S=Z$, so $D$ and $S+C$ have the same reduction. We \gitlink{audit11/decic37_jets.m}{compute} that for $t=\dm2-1$ the map $f$ of \Cref{lem:fi} is a formal immersion at $Z+\overline{c_1}$ for each of the $18$ rational cusps $c_1$. As $J_1(37)(\Q)[2]=0$, reduction is injective on $tJ_1(37)(\Q)$, so $D=S+C$ by \Cref{lem:fi}, which is a contradiction as $D$ contains no cusps.
\end{proof}

\subsection{The primes \texorpdfstring{$41$}{41} and \texorpdfstring{$43$}{43}}

\begin{proposition}\label{prop:41}
We have $41\notin S(15)$ and $41\notin S(16)$ if $Z=O_{14}$.
\end{proposition}

\begin{proof}
Here $O_{14}$ is one of the $10$ closed points of degree $14$ which are the reductions of the closed points of degree $14$ in the supports of the trace $S$ of an \gitlink{degree15/lift41_charzero.m}{explicit point} of degree $14$ on $X_1(41)$ and of its images under $\Delta(41)$, as we \gitlink{degree15/known41_lifting.m}{check}. As in the proof of \Cref{prop:37} we can assume that $\overline S=Z$. As $J_1(41)(\Q)$ is finite, the map $t=(A_3,A_5,A_7):J_1(41)\to J_1(41)^3$ kills $J_1(41)(\Q)$, and we compute that $f$ is a formal immersion at $Z+C'$ for all the $20$ (for $d=15$) and $210$ (for $d=16$, \gitlink{degree16/known41_lifting.m}{code}) possible $C'$. So $D=S+C$ by \Cref{lem:fi}, which is a contradiction.
\end{proof}

For $p=43$ and $d=16$ we need the following lemma.

\begin{lemma}\label{lem:43inj}
Reduction modulo $2$ is injective on $J_1(43)(\Q)_\tor$.
\end{lemma}

\begin{proof}
Let $Y=X_1(43)/\dm2$, where $\dm 2$ has order $7$, and let $f:X_1(43)\to Y$ be the quotient map. The curve $Y$ has genus $9$, and we \gitlink{degree16/torsion43_upper.m}{compute} that $J_Y$ contains an elliptic curve $A$ isogenous to \lmfdbec{43}{a}{1}, whose isogeny class consists of one curve with trivial rational torsion, so $A(\Q)[2]=0$, that $B=J_Y/A$ has rank $0$, and that the numbers of points of $B$ over $\F_3$, $\F_5$ and $\F_{23}$ have greatest common divisor $532=4\cdot 133$. Hence $J_Y(\Q)[2]$ injects into $B(\Q)_\tor$, which injects into $B(\F_3)$, $B(\F_5)$ and $B(\F_{23})$, so $\#J_Y(\Q)[2]\leq4$. Let $C_1,C_2,C_3$ be the rational cusps of $Y$. The divisors $38(C_2-C_1)$ and $38(C_3-C_1)$ are \gitlink{degree16/torsion43_Q_first.m}{principal}, and we \gitlink{degree16/torsion43_reduction.m}{compute} that the reductions of $x_1=19[C_2-C_1]$, $x_2=19[C_3-C_1]$ and $x_1+x_2$ are non-zero. So $J_Y(\Q)[2]=\{0,x_1,x_2,x_1+x_2\}$, and reduction is injective on it.

Now let $x\in J_1(43)(\Q)_\tor$ with $\overline x=0$. Then $2x=0$, so $f_*x\in J_Y(\Q)[2]$ has reduction $0$, and hence $f_*x=0$. We \gitlink{degree16/two43_h7.m}{compute} that $\dm 2$ acts trivially on the common kernel of $A_3$, $A_5$ and $A_7$ on $H_1(X_1(43),\Z)\otimes\F_2$, which contains $x$. Hence $0=f^*f_*x=\sum_{i=0}^6\dm2^ix=7x=x$.
\end{proof}

\begin{proposition}\label{prop:43}
We have $43\notin S(13)$ and $43\notin S(16)$ if $Z$ is one of the divisors listed for $43$ in \Cref{table:exceptions}.
\end{proposition}

\begin{proof}
The $7$ closed points $O_{12}$ are the reductions of the closed points of degree $12$ in the supports of the trace $S$ of the point of order $43$ of \Cref{sec:lower} and its images under $\Delta(43)$, as we \gitlink{degree13/mixed43_lifting.m}{check}, and the $21$ closed points $O_{14}$ are, in the same way, the reductions of the trace of an explicit point of degree $14$ on $X_1(43)$ from the tables of \cite{vanHoeij} and its images (\gitlink{degree16/lift43_14.m}{check}). For $d=13$ and $Z=O_{12}$ we \gitlink{degree13/mixed43_tangents.m}{compute} that $f$ is a formal immersion at $Z+\overline{c_1}$ for all the $21$ rational cusps $c_1$, with $t=(\dm{3}^{-1}-1)\dm{11}^{-1}A_{11}$, which kills $J_1(43)(\Q)$, and for $d=16$ and $Z=O_{14}$ we \gitlink{degree16/known43_cot.m}{compute} that $f$ is a formal immersion at all the $231$ divisors $Z+C'$, with $t$ combining several maps $\alpha A_q$, $q\in\{3,5,7\}$, with $\alpha J_1(43)(\Q)$ finite. As in the proof of \Cref{prop:37}, this is a contradiction.

The remaining cases are for $d=16$. By \Cref{lem:43inj}, $y=0$ in \Cref{lem:stab}, and $\gon_\Q X_1(43)=24$. If $Z=O_{12}$ or $Z$ is one of the two orbits $O_{15}$, then $Z$ is fixed by $h=\dm{37}$, of order $3$, and $c\in\{4,1\}$, so as $24>16+c$, \Cref{lem:stab} gives $3\mid c$, which is a contradiction. If $Z\in\{SS_{14},2SS_7,2NC_7,SS_7+NC_7\}$, then $Z$ is fixed by $h=\dm2$, of order $7$, and $c=2$, so as $24>18$ we get $7\mid2$, which is again a contradiction.
\end{proof}

\subsection{The primes \texorpdfstring{$53$}{53} and \texorpdfstring{$61$}{61}}

\begin{proposition}\label{prop:53}
We have $53\notin S(d)$ for $13\leq d\leq 21$ if $Z=SS_{13}$ or $Z=0$.
\end{proposition}

\begin{proof}
There is a unique supersingular closed point $SS_{13}$ of degree $13$, and $h=\dm4$, of order $13$, acts simply transitively on its geometric points. We \gitlink{degree13/two53_verified.json}{compute} that $J_1(53)(\Q)[2]=0$, so $y=0$ in \Cref{lem:stab}.

Suppose $Z=SS_{13}$, so $c=d-13$. If $d=13$, then as $\gon_\Q X_1(53)>13$, \Cref{lem:stab} gives $hD=D$, and by \Cref{lem:orbits} $D$ is a single orbit of $\dm4$, so its image on $X_0(53)$ is a non-cuspidal rational point, which does not exist \cite{mazur78}. If $14\leq d\leq 16$, then $d+c\leq 19<20\leq\gon_\Q X_1(53)$, and \Cref{lem:stab} gives $13\mid c$, which is a contradiction. For $17\leq d\leq 21$ we use \Cref{lem:relative} with $L=C+\pi^*Y$, where $Y=X_1(53)/\dm4$, and $t=((h-1)A_3,(h-1)A_5)$, which kills $J_1(53)(\Q)$ and $\pi^*J_Y$. We \gitlink{degree17/relative53.m}{compute} that the rank condition holds for all the possible $C'$, so $D$ lies on $L$ and contains $C\neq0$, which is a contradiction.

Suppose now that $Z=0$ and $19\leq d\leq 21$. For $d=19$ the criterion of \Cref{prop:cusps} on $X_H$ with $H$ of order $2$ fails at only $7$ multiplicity vectors, and at all the $20444$ divisors on $X_1(53)$ lying above them we \gitlink{degree19/full_cusps53.m}{compute} that $f$ is a formal immersion, with $t=((\dm{10}-1)A_3,(\dm{10}-1)A_5)$. For $d\in\{20,21\}$ we work on $X_1(53)$ with $t=((h-1)A_3,(h-1)A_5)$, and we \gitlink{degree20/full_cusps53.m}{compute} that $f$ is a formal immersion at every sum of rational cusps of degree $d$ that does not contain a full orbit of $\dm4$, while at the others the rank condition of \Cref{lem:relative} holds with $L=R+\pi^*Y$, where $R\neq0$ is the rest of the divisor. In both cases $D$ contains a cusp, which is a contradiction.
\end{proof}

\begin{proposition}\label{prop:61}
We have $61\notin S(d)$ for $d\in\{15,\dots,19,21,22\}$ if $Z\in\{SS_{15},SS_{20},O_{20}\}$.
\end{proposition}

\begin{proof}
There is a unique supersingular closed point $SS_{15}$ of degree $15$, and it is \gitlink{degree15/quotient61/verify_reduction.m}{fixed} by all of $\Delta(61)$. Let $h=\dm3=\dm2^6$, of order $5$. As the strange characters modulo $61$ have order $6$, $(h-1)J_1(61)(\Q)$ is finite, and we \gitlink{degree15/quotient61/verify_sage.py}{compute} that $J_1(61)(\Q)[2]=0$, so $y=0$ in \Cref{lem:stab}. Note that $\gon_\Q X_1(61)>18$.

Suppose $Z=SS_{15}$. If $d=16$, then $c=1$, and \Cref{lem:stab} gives $5\mid 1$. If $d\in\{17,18,19,21,22\}$, we use \Cref{lem:relative} with $L=C+\pi^*Y^{(3)}$, where $Y=X_1(61)/\dm h$, and $t=((h-1)A_3,(h-1)A_5)$, and we \gitlink{degree17/relative61.m}{compute} that the rank condition holds for all possible $C'$, so $D$ contains $C\neq0$, which is impossible. In the same way, for $d\in\{21,22\}$ and $Z\in\{SS_{20},O_{20}\}$, where $O_{20}$ is an orbit of $6$ closed points which are pullbacks from $Y$, we use $L=C+\pi^*Y^{(4)}$ (\gitlink{degree21/relative61_ss20.m}{$SS_{20}$}, \gitlink{degree21/relative61_o20.m}{$O_{20}$}).

It remains to deal with $d=15$ and $Z=SS_{15}$, so $c=0$. By \Cref{lem:stab} $hD=D$, so by \Cref{lem:orbits} $D$ is a pullback from $Y$, $D=\pi^*D_Y$, and $D_Y$ is a closed point of degree $3$. Let $W=X_1(61)/\dm4$, where $\dm4$ has order $15$. Then $W$ has genus $8$, the image $s$ of $\dm2$ is an involution of $W$, and $W/\dm s=X_0(61)$. The image $D_W$ of $D_Y$ on $W$ is a closed point of degree $3$, as a rational point would map to a non-cuspidal rational point on $X_0(61)$. Its reduction is fixed by $s$. The Jacobian $J_W$ corresponds to the characters modulo $61$ of order at most $2$, so $(s-1)J_W(\Q)$ is finite by Kato's theorem as in \Cref{lem:finite}, and $J_W(\Q)[2]=0$ since $J_1(61)(\Q)[2]=0$ and the pullback to $X_1(61)$ is injective on $2$-torsion (its composition with the pushforward is multiplication by $15$). Hence, as in \Cref{lem:stab}, $sD_W\sim D_W$, and as $\gon_\Q W\geq 4$ by \Cref{lem:W61} below, $sD_W=D_W$. As $2$ and $3$ are coprime, $s$ fixes each of the three geometric points of $D_W$. If $w$ is one of them and $x$ is a point of $D$ above it, then $\dm2x=gx$ for some $g\in\dm4$, so $x$ is fixed by $\dm2g^{-1}\neq1$. This contradicts \Cref{lem:CM}, as $x$ has degree $15<20$.
\end{proof}

\begin{lemma}\label{lem:W61}
The curve $W=X_1(61)/\dm4$ has no map of degree $2$ or $3$ to $\PP^1$ defined over $\Q$.
\end{lemma}

\begin{proof}
The curve $X_0(61)$ has genus $4$ and no map of degree at most $3$ to $\PP^1$ over $\Q$: its canonical model is a complete intersection of a quadric and a cubic in $\PP^3$, and the discriminant of the quadric is not a square, so the two $g^1_3$'s are conjugate over $\Q(i)$ (see \cite{HS,jeon} and \gitlink{degree15/quotient61/verify_sage.py}{our check}). Suppose that $u:W\to\PP^1$ has degree $e\in\{2,3\}$. As $g(W)=8>(e-1)^2$, the Castelnuovo-Severi inequality implies that $u$ is the unique map of degree $e$ to $\PP^1$ up to automorphisms of $\PP^1$, so $u\circ s=\beta\circ u$ for some $\beta\in\mathrm{PGL}_2(\Q)$ with $\beta^2=1$. If $\beta=1$, then $u$ factors through $W\to X_0(61)$, which is impossible as $e\leq 3$. Otherwise $u$ induces a map $X_0(61)\to\PP^1/\dm\beta\simeq\PP^1$ of degree $e$, which is again impossible.
\end{proof}

\subsection{The prime \texorpdfstring{$73$}{73}}

\begin{proposition}\label{prop:73}
We have $73\notin S(d)$ for $9\leq d\leq 22$ if $Z$ is a sum of points $SS_6$, $NC_9$ and $SS_{18}$.
\end{proposition}

\begin{proof}
Let $h=\dm2$, of order $9$. It fixes each of the $4$ supersingular closed points of degree $6$ and the $4$ non-rational cusps (of degree $9$), and the points of degree $18$, and the subgroup $\dm8$ of order $3$ fixes the geometric points of degree $6$. We \gitlink{audit11/rational_two.m}{compute} that $h$ acts trivially on the common kernel of $A_3$ and $A_5$ on $H_1(X_1(73),\Z)\otimes\F_2$, which contains $J_1(73)(\Q)[2]$, so $y=0$ in \Cref{lem:stab} ii), and $\gon_\Q X_1(73)>26$. As $d<24=(73-1)/3$, the orbits of $\dm 2$ on the support of $D$ have length $9$ (\Cref{lem:orbits}). Let $e=\deg Z$, so $c=d-e$ and $e\in\{6,9,12,15,18,21\}$.

Suppose first that $d+c=2d-e\leq 26$. Then \Cref{lem:stab} gives $hD=D$ and $9\mid c$, and as the orbits of $h$ have length $9$, $9$ divides $d$. Hence $d\in\{9,18\}$ and $D$ is a pullback from $Y=X_1(73)/\dm2$, of a divisor $D'$ of degree $d/9$. If $d=9$, then the image of $D'$ on $X_0(73)$ is a non-cuspidal rational point, which does not exist \cite{mazur78}. If $d=18$, then $D'$ maps to a non-cuspidal rational or quadratic point on $X_0(73)$. These were determined by Box \cite[Section 4.7]{Box}, and they have $11$ different $j$-invariants ($3$ of them quadratic and not CM, and $8$ rational CM $j$-invariants). For each of them we \gitlink{degree18/quotient_quadratic.m}{check}, using the action of Frobenius on $E[73]$ at a prime which splits in the corresponding quadratic field, that no point of $Y$ of degree at most $2$ lies above it.

Suppose now that $2d-e>26$. Then $d\geq 17$, and $Z$ and $d$ are as in \Cref{table:73}. In these cases we use \Cref{lem:relative} with $t=((h-1)A_3,(h-1)A_5)$ and
$$L=Z_0+R+\pi^*Y^{(r)},$$
where $Z_0$ is the divisor of degree $6k$ whose points are the unique lifts of the geometric points of $SS_6$ in $Z$ that are fixed by $\dm8$ (here $k$ is the number of points $SS_6$ in $Z$), $C'=\overline{R}+\pi^*E'$ with $E'$ a sum of $r'$ cusps of $Y$, $R$ has no full orbit of $\dm2$ in its support, and $r=r'$ plus the number of points $NC_9$ in $Z$ (which are pullbacks of cusps of $Y$). We \gitlink{degree17/relative73.m}{compute} (\gitlink{degree18/relative73_nc.m}{$NC_9$}, \gitlink{degree21/relative73_multiple.m}{two or three non-ordinary points}) that the rank condition holds in all cases, so $D$ lies on $L$. If $R\neq0$, then $D$ contains a cusp, which is impossible. If $R=0$ and $k>0$, then $D$ contains a point fixed by $\dm8$, which contradicts \Cref{lem:CM}. If $R=0$ and $k=0$, then $Z=NC_9$, $d=18$, and $D$ is a pullback from $Y$, which we have already excluded.
\end{proof}

\begin{table}[H]
\centering
\begin{tabular}{rl}
\toprule
$d$ & $Z$\\
\midrule
$17$ & $SS_6$\\
$18$ & $SS_6$, $NC_9$\\
$19$ & $SS_6$, $NC_9$\\
$20$ & $SS_6$, $NC_9$, $2SS_6$\\
$21$ & $SS_6$, $NC_9$, $2SS_6$, $SS_6+NC_9$\\
$22$ & $SS_6$, $NC_9$, $2SS_6$, $SS_6+NC_9$\\
\bottomrule
\end{tabular}
\medskip
\caption{The reductions at $73$ for which we use \Cref{lem:relative}.}
\label{table:73}
\end{table}

\subsection{The primes \texorpdfstring{$89$}{89} and \texorpdfstring{$97$}{97}}

\begin{proposition}\label{prop:89}
We have $89\notin S(d)$ for $11\leq d\leq 22$ if $Z\in\{NC_{11},2NC_{11},SS_{22}\}$.
\end{proposition}

\begin{proof}
Let $h=\dm2$, of order $11$, which fixes the $4$ non-rational cusps. We \gitlink{degree16/certify_anti89.m}{compute} that $A_3$ has odd degree on $(h-1)J_1(89)$, so $y=0$ in \Cref{lem:stab} iii), and $\gon_\Q X_1(89)>39$. If $Z=NC_{11}$, then $c=d-11$ and $d+c\leq 33$, so \Cref{lem:stab} gives $hD=D$ and $11\mid c$. If $c=0$, then $d=11$, $D$ is a single orbit of $\dm 2$ by \Cref{lem:orbits}, and its image on $X_0(89)$ is a non-cuspidal rational point, which does not exist. So $c=11$ and $d=22$.

It remains to deal with $d=22$ and $Z\in\{NC_{11},2NC_{11},SS_{22}\}$. The supersingular closed points of degree $22$ form $5$ orbits under $\Delta(89)$, of lengths $4,4,4,2,1$. For the first three, $\alpha A_3Z$ or $\alpha A_5Z$ is not principal (\Cref{prop:sieve}), and the other two are fixed by $h$. Hence, by \Cref{lem:stab} and \Cref{lem:orbits}, in all the remaining cases $D$ is a pullback from $Y=X_1(89)/\dm2$, of a divisor $D'$ of degree $2$. The curve $Y$ has genus $27$, and we \gitlink{degree22/check_quotient89.py}{compute} that $\#Y(\F_2)=8$ and $\#Y(\F_4)=38$, so $\gon_\Q Y\geq\gon_{\F_4}Y\geq8$.

If $Z$ contains a non-rational cusp, then $\overline{D'}$ is a sum of two cusps of $Y_{\F_2}$, which are defined over $\F_2$. We \gitlink{degree22/quotient89_cusps.m}{compute} that the map $Y^{(2)}\to J_Y$ given by $t=(h'-1)(T_q-\dm q-q)A_q$, where $h'$ generates $\Delta(89)/\dm2$, is a formal immersion at all the $36$ sums of two such cusps. As $t$ kills $J_Y(\Q)$ and all the classes of differences of cusps, $t$ maps $D'$ and the cuspidal divisor lifting $\overline{D'}$ (which is defined over an unramified extension of $\Q_2$) to $0$, so they are equal by \cite[Lemma 4.2]{DKSS}, which is a contradiction.

If $Z=SS_{22}$, let $\iota$ be the element of order $2$ of $\Delta(89)/\dm2$, which is cyclic of order $4$. It fixes $\overline{D'}$, $(\iota-1)J_Y(\Q)$ is finite, and so, as in \Cref{lem:stab}, $2\iota D'\sim 2D'$. As $\gon_\Q Y>4$, it follows that $\iota D'=D'$. A point of $Y$ fixed by $\iota$ lies below a point of $X_1(89)$ fixed by a non-trivial element of $\Delta(89)$, which has degree at least $30$ by \Cref{lem:CM}, so it has degree at least $3$. Hence $\iota$ has no fixed points in the support of $D'$, so $D'=y+\iota(y)$ for a point $y$ of degree at most $2$, and the image of $D'$ on $Y/\dm\iota$ is twice a rational point. Its image on $X_0(89)$ is a non-cuspidal rational point, which does not exist.
\end{proof}

\begin{proposition}\label{prop:97}
We have $97\notin S(d)$ for $16\leq d\leq 22$ if $Z=SS_{16}$.
\end{proposition}

\begin{proof}
The two supersingular closed points of degree $16$ are fixed by the subgroup $H_4\subseteq\Delta(97)$ of order $4$, and their geometric points are fixed by the subgroup $H_3$ of order $3$. As the strange characters modulo $97$ have order $12$, $(h-1)J_1(97)(\Q)$ is finite for $h\in H_4$, and \Cref{prop:gonality} gives $\gon_\Q X_1(97)\geq49$.

Let $h$ be a generator of $H_4$. If $16\leq d\leq 19$, then $c=d-16\leq 3$ and $2(d+c)\leq 44<49$, so \Cref{lem:stab} gives $hD=D$ and $4\mid c$, so $c=0$ and $d=16$. As $16<(97-1)/3$, $D$ is a pullback from $X_1(97)/H_4$ (\Cref{lem:orbits}). Let $H_8$ be the subgroup of order $8$, $W=X_1(97)/H_8$, of genus $41$, and let $h'$ be the element of order $3$ of $\Delta(97)/H_8$. The image $D_W$ of $D$ on $W$ is a closed point of degree $m\in\{2,4\}$, and its reduction is fixed by $h'$. The group $(h'-1)J_W(\Q)$ is finite, as $J_W$ corresponds to the characters of order dividing $6$, and we \gitlink{degree16/certify_anti97.m}{compute} that $A_3$ has odd degree on $(h'-1)J_W$. As in \Cref{lem:stab}, it follows that $h'D_W\sim D_W$, and as $\gon_\Q W\geq 49/8$ (a map $W\to\PP^1$ of degree $k$ gives a map $X_1(97)\to\PP^1$ of degree $8k$), $h'D_W=D_W$. As $3\nmid m$, $h'$ fixes each geometric point of $D_W$, so the points of $D$ are fixed by non-trivial elements of $\Delta(97)$, which contradicts \Cref{lem:CM}.

If $20\leq d\leq 22$, we use \Cref{lem:relative} with $Y=X_1(97)/H_3$ and $L=Z_0+R+\pi^*Y^{(r)}$, where $Z_0$ is the divisor of degree $16$ whose points are the lifts of the geometric points of $SS_{16}$ fixed by $H_3$, and $R$ and $r$ are as in the proof of \Cref{prop:73}. Here we use $t=(h_3-1)\beta_nA_q$, where $h_3$ generates $H_3$, $q$ is odd, and $\beta_n=T_n-a_n$, where $a_n\in\Z[\Delta(97)]$ is such that $\beta_n$ kills the abelian subvariety of $J_1(97)$ corresponding to the strange characters (on which $T_n$ acts through $\Z[\Delta(97)]$, as this subvariety has dimension $1$ over the field of values of the characters). Then $t$ kills $J_1(97)(\Q)$ and $\pi^*J_Y$, and we \gitlink{degree20/sharp97.m}{compute} that the rank condition holds. So $D$ contains $R$, which is impossible if $R\neq0$, or $Z_0$, which contradicts \Cref{lem:CM}.
\end{proof}

\subsection{Reductions supported on the rational cusps at \texorpdfstring{$101$}{101}, \texorpdfstring{$103$}{103} and \texorpdfstring{$131$}{131}}

\begin{proposition}\label{prop:cusps2}
Let $(d,p)$ be $(20,101)$, $(21,101)$, $(22,101)$, or $p=103$ and $17\leq d\leq 22$, or $p=131$ and $18\leq d\leq 22$. Then $Z\neq 0$.
\end{proposition}

\begin{proof}
For $p=101$ we use $t=(h-1)\beta_nA_q$ with $h$ a generator of $\Delta(101)$ and $\beta_n$ as in the proof of \Cref{prop:97}, for $n\in\{2,3,5,7\}$, and we \gitlink{degree20/sharp101.m}{compute} that $f$ is a formal immersion at all sums of rational cusps of degree $d$ that do not contain all the $50$ rational cusps, which is all of them.

Let $p=103$ (resp.\ $p=131$), let $H\subseteq\Delta(p)$ be the subgroup of order $3$ (resp.\ $5$), and $Y=X_H$. Then $Y\to X_0(p)$ has degree $17$ (resp.\ $13$) and the rational cusps of $Y$ form one fibre $F$ of this map. On $Y$, the criterion of \Cref{prop:cusps} holds at all sums of rational cusps of degree $d$ that do not contain $F$. If $\overline{\pi_*D}$ contains $F$, write $\overline{\pi_*D}=F+\overline R$. We use \Cref{lem:relative} on $Y^{(d)}$ with $L=R+\rho^*X_0(p)$, where $\rho:Y\to X_0(p)$, and $t=(\dm a-1)(T_3-\dm 3-3)$ on $J_Y$, where $\dm a$ generates $\Delta(p)/H$, which kills $J_Y(\Q)$ (as $p$ is not strange) and $\rho^*J_0(p)$. We \gitlink{degree19/verify_relative_cusps_intrinsic.py}{compute} that the rank condition holds, so $\pi_*D$ lies on $L$. If $R\neq0$, then $\pi_*D$ contains a cusp, which is impossible. If $R=0$, then $p=103$, $d=17$, and $\pi_*D$ is the pullback of a non-cuspidal rational point on $X_0(103)$, which does not exist.
\end{proof}

\subsection{The primes \texorpdfstring{$109$}{109}, \texorpdfstring{$113$}{113}, \texorpdfstring{$127$}{127} and \texorpdfstring{$137$}{137}}
At these primes we only use $2y=0$ in \Cref{lem:stab}, so we need $\gon_\Q X_1(p)>2(d+c)$, and $d<(p-1)/3$, so $D$ is a pullback from $X_H$ whenever it is $H$-invariant.

\begin{proposition}\label{prop:109}
We have $109\notin S(d)$ for $12\leq d\leq 22$ if $Z\in\{SS_9,NC_{18},2SS_9,SS_9+SS_{12}\}$.
\end{proposition}

\begin{proof}
All of these are fixed by $h=\dm{23}$, of order $18$, and $\gon_\Q X_1(109)>58$. If $Z=SS_9$ and $d\leq 19$, or if $Z$ is one of the others, then $2(d+c)\leq58$, so $hD=D$ and $18\mid c$. As $c\leq13$, this gives $c=0$, so $Z=NC_{18}$ or $2SS_9$ with $d=18$, or $Z=SS_9+SS_{12}$ with $d=21$. In the last case $18\nmid 21$, which contradicts \Cref{lem:orbits}, and in the first two $D$ is a single orbit of $\dm{23}$, whose image on $X_0(109)$ is a non-cuspidal rational point. If $Z=SS_9$ and $d\geq 20$, we use \Cref{lem:relative} with $L=C+\pi^*Y$ or $L=R+\pi^*Y^{(2)}$, where $Y=X_1(109)/\dm{h^2}$ and $\dm{h^2}$, of order $9$, acts freely on the geometric points of $SS_9$, and $t=(h^2-1)A_q$ (\gitlink{degree20/relative109.m}{code}). As before, $D$ contains a cusp, which is a contradiction.
\end{proof}

\begin{proposition}\label{prop:113}
We have $113\notin S(d)$ for $13\leq d\leq 22$ if $Z$ is one of the divisors listed for $113$ in \Cref{table:exceptions}.
\end{proposition}

\begin{proof}
Here $O_8$ denotes one of the $7$ ordinary closed points of degree $8$ which are fixed by $h_8=\dm{40}$, of order $8$, and we have $\gon_\Q X_1(113)>63$. The element $h_2=\dm2$, of order $14$, fixes the closed points $SS_7$ and $NC_{14}$, and the involution $\iota=\dm{98}$ fixes the points $SS_7$ and $O_8$.

If $Z\in\{2SS_7,3SS_7,NC_{14},SS_7+NC_{14}\}$, then $2(d+c)\leq 60$, so $h_2D=D$ and $14\mid c$, so $c=0$. Then $d\in\{14,21\}$, and as $14\nmid21$, $d=14$ and $D$ is a single orbit of $h_2$, whose image on $X_0(113)$ is a non-cuspidal rational point, which is impossible.

If $Z=O_8$ and $d\leq19$, or $Z=2O_8$, then $2(d+c)\leq 62$, so $h_8D=D$ and $8\mid c$. Hence $d=16$ and $c=0$ or $c=8$, and $D$ is a pullback of a divisor of degree $2$ from $X_1(113)/\dm{h_8}$. Its image on $X_0(113)$ is a non-cuspidal point of degree at most $2$, which has complex multiplication, as all the quadratic points on $X_0(113)$ are CM points \cite[Table 14]{AKMNOV} and there are no non-cuspidal rational points. But then by Silverberg's theorem (see the proof of \Cref{lem:CM}), a point of order $113$ on $E$ is defined over a field of degree at least $112/6>16$, which is a contradiction. If $Z=O_8$ and $d\geq 20$, we use \Cref{lem:relative} with $L=R+\pi^*Y^{(r)}$, where $\pi$ is the quotient by $\dm{h_8}$ near $Z$, which is \'etale there, and $t=((h_8-1)A_2,(h_8-1)A_3)$ (\gitlink{degree20/relative113_ordinary.m}{code}).

If $Z=O_8+SS_7$, then $2(d+c)\leq 58$, so $\iota D=D$ and $2\mid c=d-15$. So $d$ is odd, and the orbits of $\iota$ on $D$ have length $2$ by \Cref{lem:orbits}, which is a contradiction. Finally, if $Z=O_8+NC_{14}$ and $d=22$, we use \Cref{lem:relative} with $t=(h_8-1)A_qA_{q'}$, and $L=Z_1+\pi^*Y$, where $Z_1$ is the lift of $NC_{14}$ to a cuspidal divisor (\gitlink{degree22/relative113_nc.m}{code}). Hence $D$ contains a cusp, which is a contradiction.
\end{proof}

\begin{proposition}\label{prop:127}
We have $127\notin S(d)$ for $17\leq d\leq 22$ if $Z$ is one of the divisors listed for $127$ in \Cref{table:exceptions}.
\end{proposition}

\begin{proof}
The element $h_7=\dm2$, of order $7$, fixes the $9$ non-rational cusps and the $48$ supersingular closed points of degree $14$, and the ordinary closed points $O_9$ are fixed by the subgroup $H_9$ of order $9$. We have $\gon_\Q X_1(127)>79$, and $A_2$ kills $J_1(127)(\Q)_\tor$ by \Cref{lem:torsion} b).

Suppose that $Z$ is a sum of points $NC_7$ and $SS_{14}$. Then $2(d+c)\leq 74$, so $h_7D=D$ and $7\mid c$. This is impossible unless $d=21$, in which case $D$ is a pullback from $Y=X_1(127)/\dm2$, of a divisor $D'$ of degree $3$. The curve $Y$ has genus $88$, and its points over $\F_2$ which are images of cusps are the $9$ rational cusps and the $9$ images of the non-rational cusps. If $Z$ consists of non-rational cusps, then $\overline{D'}$ is a sum of three of these $18$ points, and we \gitlink{degree21/quotient127_cusps.m}{compute} that the map $Y^{(3)}\to J_Y$ given by $t=(h'-1)(T_q-\dm q-q)A_q$, $q\in\{2,3\}$, is a formal immersion at all the $1140$ such sums, which shows, as in the proof of \Cref{prop:89}, that $D'$ is cuspidal. If $Z=SS_{14}+C_7$, then for $5$ of the $6$ orbits of points $SS_{14}$, $\alpha A_qZ$ is not principal for some $q\in\{3,5,7\}$, and for the remaining orbit (of length $3$), $\overline{D'}=S+\overline{c_1}$, where $c_1$ is a rational cusp of $Y$ and $S$ is a supersingular point of degree $2$ fixed by the element $h_3'$ of order $3$ of $\Delta(127)/\dm2$. We use \Cref{lem:relative} on $Y^{(3)}$ with $L=S_0+c_1$, where $S_0$ is the unique lift of $S$ fixed by $h_3'$, and $t=(h_3'-1)(T_q-\dm q-q)$ (\gitlink{degree21/quotient127_ss.m}{code}), so $D'$ contains a cusp. The case $SS_{14}+NC_7$ is excluded by \Cref{prop:sieve}.

If $Z=O_9$ or $Z=2O_9$, then $2(d+c)\leq 72$, so $D$ is $H_9$-invariant and $9\mid c$. Hence $d=18$ and $D$ is a pullback of a divisor of degree $2$ from $X_1(127)/H_9$. Its image on $X_0(127)$ is a non-cuspidal point of degree at most $2$, which is a CM point by \cite[Table 16]{AKMNOV}, and as in the proof of \Cref{prop:113} this contradicts Silverberg's theorem, as $126/6=21>18$.

Finally, if $Z=O_9+NC_7$, we use \Cref{lem:relative} with $L=C+\pi_9^*(X_1(127)/H_9)+\pi_7^*(X_1(127)/\dm2)$, which is smooth of dimension $2$ as $H_9$ and $\dm2$ act freely on the supports of $O_9$ and $NC_7$, and $t=(h_9-1)(h_7-1)A_q$, $q\in\{2,3\}$, where $h_9$ generates $H_9$ (\gitlink{degree17/relative127.m}{code}). So $D$ contains $C\neq0$, which is a contradiction.
\end{proof}

\begin{proposition}\label{prop:137}
We have $137\notin S(d)$ for $19\leq d\leq 22$ if $Z\in\{SS_{17},O_{17}\}$.
\end{proposition}

\begin{proof}
The point $SS_{17}$ is fixed by $\dm2$, of order $34$, and $O_{17}$ is \gitlink{degree19/stabilizer137_ordinary.m}{fixed} by $\dm{81}$, of order $17$. In both cases $c=d-17\in\{2,3,4,5\}$ and $2(d+c)\leq 54<\gon_\Q X_1(137)$, so \Cref{lem:stab} gives $34\mid c$ or $17\mid c$, which is a contradiction.
\end{proof}

\subsection{Proof of \texorpdfstring{\Cref{thm:main}}{Theorem 1.1}}
By \Cref{prop:lower} all the primes listed in \Cref{thm:main} are in $S(d)$. Conversely, let $9\leq d\leq 22$ and let $p\in S(d)$ be a prime which is not listed. By \Cref{prop:large,prop:medium}, $p$ is one of the primes of \Cref{table:local}. By \Cref{prop:cusps,prop:sieve,prop:cusps2}, the reduction of the trace divisor is one of those listed in \Cref{table:exceptions}, and these are excluded by \Cref{prop:37,prop:41,prop:43,prop:53,prop:61,prop:73,prop:89,prop:97,prop:109,prop:113,prop:127,prop:137}.
